% Irodov Part IV §4.1 original 51--65, source PDF pp.169--172
% printed pp.160--163, figures 4.15--4.23 checked.
\oscproblem{51}{160}
均匀细杆长度 $\ell$，绕穿过杆上某点并垂直于杆的水平轴作小振动。
求悬轴距质心多远时周期最短，并求最小周期。
\exsolution
设轴距质心为 $x>0$，
对轴转动惯量
$J=m(\ell^2/12+x^2)$，
重力回复力矩系数 $K=mgx$。
故
\[
T(x)=2\pi\sqrt{\frac{\ell^2/12+x^2}{gx}}
=2\pi\sqrt{\frac1g\left(x+\frac{\ell^2}{12x}\right)}.
\]
极小点满足
$\mathrm d(x+\ell^2/(12x))/\mathrm dx
=1-\ell^2/(12x^2)=0$，于是
\[
\boxed{x_{\rm opt}=\frac{\ell}{2\sqrt3}},
\qquad
\boxed{T_{\min}=2\pi\sqrt{\frac{\ell}{g\sqrt3}}}.
\]
\excheck 二阶导数
$\ell^2/(6x^3)>0$，
确认是最小而非最大。

\oscproblem{52}{160}
边长相等的均匀薄等边三角形板，
其高度为 $h$。
以其中一条边为水平悬挂轴作小角平面摆动。
求周期和等效摆长。
\exsolution
三角形板的质心距底边
$d=h/3$。
对底边轴转动惯量
$J=m h^2/6$
（可按与底边平行的薄条积分）。
因此
\[
\omega_0^2=\frac{mgd}{J}
=\frac{2g}{h},
\qquad
\boxed{T=2\pi\sqrt{\frac{h}{2g}}
=\pi\sqrt{\frac{2h}{g}}}.
\]
与普通单摆周期
$T=2\pi\sqrt{\ell_{\rm red}/g}$
比较：
\[
\boxed{\ell_{\rm red}=\frac{J}{md}=\frac h2}.
\]
\excheck 对悬轴的转动惯量应为 $mh^2/6$，
不是三角板面内垂直板面的转动惯量。

\oscproblem{53}{161}
图4.15：水平光滑圆盘绕竖直轴 O
以恒角速度 $\Omega$ 转动。
均匀细杆 AB 长 $\ell$，
绕盘上距轴 O 为 $a$ 的竖直轴 A 自由小角摆动。
求相对圆盘径向稳定方向的小振动角频率。
\exsolution
在与圆盘共同转动的参考系，
杆上一质元到 O 的距离平方
$r^2=a^2+2as\cos\theta+s^2$，
其中 $0\le s\le\ell$ 从 A 沿杆测距。
离心势能
$U_{\rm cen}=-\tfrac12\Omega^2\int r^2\,\mathrm dm$，
与 $\theta$ 有关部分
$U_{\rm cen}=-\tfrac12m\Omega^2a\ell\cos\theta$。
小角展开得回复刚度
$K=m\Omega^2a\ell/2$。
杆对 A 的转动惯量
$J=m\ell^2/3$，
因此
\[
\boxed{\omega_0=
\sqrt{\frac KJ}
=\Omega\sqrt{\frac{3a}{2\ell}}}.
\]
\excheck 角频率维度为 $\mathrm{s^{-1}}$，
故不能只写 $\frac{3a\Omega^2}{2\ell}$
（那是角频率的平方）。

\oscproblem{54}{161}
图4.16：半径为 $R$、对轴转动惯量为 $J$ 的滑轮绕轴无摩擦，
一侧悬质量 $m$，另一侧接轻弹簧刚度 $k$。
细线不滑动且无质量。
求平衡位置附近的小振动角频率。
\exsolution
滑轮角位移 $\theta$ 对应悬块及弹簧端点的线位移
$x=R\theta$。
动能
$T=\tfrac12m\dot x^2+
\tfrac12J(\dot x/R)^2
=\tfrac12(m+J/R^2)\dot x^2$。
重力只确定静平衡位置，
微小变形的势能为
$\tfrac12kx^2$。
故
\[
\boxed{\omega_0=
\sqrt{\frac{k}{m+J/R^2}}}.
\]
\excheck 滑轮惯性体现为附加等效质量 $J/R^2$，
不能忽略。

\oscproblem{55}{161}
图4.17：质量为 $M$、半径为 $R$ 的均匀圆柱滑轮，
轴 O 固定且无摩擦；
右侧悬线挂重物 A，
滑轮边缘另一处固定点质量 $m$。
在图示角 $\alpha$ 处静平衡，求小振动角频率。
\exsolution
设悬挂重物 A 的质量为 $m_A$。
平衡时悬重力矩
$m_AgR$ 与边缘点质量的力矩
$mgR\sin\alpha$ 抵消，所以
$m_A=m\sin\alpha$。
小角偏离 $\delta\theta$ 时，
仅边缘点质量的重力矩发生线性回复变化，
刚度为
$K=mgR\cos\alpha$。
总转动惯量
\[
J_{\rm tot}=\frac12MR^2+mR^2+m_AR^2
=R^2\left[\frac M2+m(1+\sin\alpha)\right].
\]
所以
\[
\boxed{\omega_0=
\sqrt{\frac{2mg\cos\alpha}
{R[M+2m(1+\sin\alpha)]}}}.
\]
\excheck 右侧悬重物虽未单独给质量，
但可由平衡条件消去；其平动惯性不可忽略。

\oscproblem{56}{161}
半径为 $r$ 的均匀实心圆柱，
沿半径 $R>r$ 的固定圆柱\emph{内壁}
无滑动滚动，在最低位置附近作小振动。
求周期。
\exsolution
移动圆柱的质心沿半径
$\rho=R-r$ 的圆弧运动。
设质心相对最低点转角为 $\theta$。
无滑动约束使自身角速度模
$\omega_s=\rho\dot\theta/r$。
动能
\[
T_k=\tfrac12m\rho^2\dot\theta^2+
\tfrac12\left(\frac12mr^2\right)
\left(\frac{\rho\dot\theta}{r}\right)^2
=\frac34m\rho^2\dot\theta^2.
\]
势能
$U=mg\rho(1-\cos\theta)
\simeq\tfrac12mg\rho\theta^2$。
故
$\omega_0^2=2g/(3\rho)$，
\[
\boxed{T=2\pi\sqrt{\frac{3(R-r)}{2g}}}.
\]
\excheck 内壁几何决定质心圆轨道半径为 $R-r$，
不能用 $R+r$。

\oscproblem{57}{161}
图4.18：质量为 $m$ 的均匀实心圆柱在水平面上无滑动滚动。
其顶部由左右弹簧牵引，合刚度为 $k$。
求小振动周期。
\exsolution
令圆柱质心水平位移为 $x$，无滑动时转角
$\theta=x/R$。
顶部连接点水平位移
$x+R\theta\simeq2x$，
两弹簧储能
$U=\tfrac12k(2x)^2=2kx^2$，
等效水平刚度 $k_x=4k$。
平动和自转动能
\[
T_k=\frac12m\dot x^2+
\frac12\left(\frac12mR^2\right)
\frac{\dot x^2}{R^2}
=\frac34m\dot x^2,
\]
等效质量 $m_{\rm eff}=3m/2$。
因此
\[
\boxed{T=2\pi\sqrt{\frac{3m/2}{4k}}
=\pi\sqrt{\frac{3m}{2k}}}.
\]
\excheck 图中弹簧连接于圆柱顶端而不是圆心，
因此弹簧端位移约是质心位移的两倍。

\oscproblem{58}{161}
两个质量分别为 $m_1,m_2$ 的立方体，
用刚度 $k$ 的轻弹簧相连，
置于水平光滑面上，
压缩弹簧后同时松手。
求内部相对运动的固有角频率。
\exsolution
以两物体坐标 $x_1,x_2$ 描述，
相对坐标
$\xi=x_2-x_1-\ell_0$，
则
$m_1\ddot x_1=k\xi$、
$m_2\ddot x_2=-k\xi$。
相减：
$\ddot\xi+k(1/m_1+1/m_2)\xi=0$。
定义约化质量
$\mu=m_1m_2/(m_1+m_2)$，
所以
\[
\boxed{\omega_0=\sqrt{k\left(\frac1{m_1}+
\frac1{m_2}\right)}=\sqrt{\frac k\mu}}.
\]
\excheck 质心作匀速运动与弹簧内振动解耦。

\oscproblem{59}{162}
图4.19：两滑球质量
$m_1=1.0$kg、$m_2=2.0$kg，
由刚度 $k=24$N/m 的轻弹簧相连，
在光滑水平杆上滑动。
起初弹簧未伸长，突然给予左球
$v_1=12$cm/s 初速。
求(a)内部振动角频率；
(b)振动能及振幅。
\exsolution
约化质量
$\mu=m_1m_2/(m_1+m_2)=2/3$kg，
相对坐标满足
$\mu\ddot\xi+k\xi=0$，
所以
\[
\boxed{\omega_0=\sqrt{k/\mu}
=\sqrt{24/(2/3)}=6.0\,\mathrm{s^{-1}}}.
\]
起初相对位移 $\xi(0)=0$，
相对速度的模 $|\dot\xi(0)|=v_1$，
振幅
$\boxed{A=v_1/\omega_0=0.020\,\mathrm m=2.0\,\mathrm{cm}}$。
振动能
\[
\boxed{E_{\rm osc}=\frac12\mu v_1^2
\simeq4.8\,\mathrm{mJ}\approx5\,\mathrm{mJ}}.
\]
\excheck 总系统还有质心平动能；
上述 $E_{\rm osc}$ 仅为内部弹簧振动能。

\oscproblem{60}{162}
两个圆盘的轴向转动惯量分别为
$J_1,J_2$，由扭转系数为 $k$ 的细杆连接。
求二者相对小扭转运动周期。
\exsolution
设两圆盘角位移分别
$\theta_1,\theta_2$。
扭转势能
$U=\tfrac12k(\theta_1-\theta_2)^2$。
角动量整体守恒，可用相对角
$\phi=\theta_1-\theta_2$ 描述，
对应约化转动惯量
$J_{\rm red}=J_1J_2/(J_1+J_2)$。
于是
$J_{\rm red}\ddot\phi+k\phi=0$，
\[
\boxed{T=2\pi\sqrt{\frac{J_1J_2}{k(J_1+J_2)}}}.
\]
\excheck 固有整体刚体转动为零频模式；
这里给的是两盘\emph{相对}扭转振动周期。

\oscproblem{61}{162}
图4.20：用两个相同轻弹簧模拟线性
O—C—O 的二氧化碳分子。
三颗小球质量分别为
$m_O,m_C,m_O$，
可按图示两种模态振动。
求两种振动角频率之比。
\exsolution
设每根弹簧刚度 $k$。
对称振型：两端 O 同时向中心或背离中心等幅运动，
中央 C 静止，单端满足
$m_O\ddot x=-kx$，
$\omega_1^2=k/m_O$。
另一振型两 O 同向等幅位移 $x$，
C 反向位移 $y$，总质心守恒
$2m_Ox+m_Cy=0$。
端球的弹力为
$-k(x-y)=-k(1+2m_O/m_C)x$，
故
\[
\omega_2^2=\frac{k}{m_O}
\left(1+\frac{2m_O}{m_C}\right).
\]
因此
\[
\boxed{\frac{\omega_2}{\omega_1}
=\sqrt{1+\frac{2m_O}{m_C}}
=\sqrt{1+\frac{32}{12}}\simeq1.9}.
\]
\excheck 模态顺序依据原版图4.20的箭头；
振动比值与弹簧刚度无关。

\oscproblem{62}{162}
图4.21：充有理想气体的封闭圆筒，
活塞截面积为 $S$，质量为 $m$，
静平衡时将筒内两侧均分成体积 $V_0$，
两侧气压为 $p_0$。
活塞微小位移后释放，
气体变化为绝热过程，忽略摩擦。
求振动角频率。
\exsolution
活塞位移 $x$ 时，
左、右气体体积分别
$V_0+Sx,V_0-Sx$。
由绝热关系 $pV^\gamma=\mathrm{const}$，
一阶展开：
\[
p_L\simeq p_0\left(1-\frac{\gamma Sx}{V_0}\right),
\quad
p_R\simeq p_0\left(1+\frac{\gamma Sx}{V_0}\right).
\]
压力差给回复力
$F=S(p_L-p_R)
=-2\gamma p_0S^2x/V_0$，
所以
\[
\boxed{\omega_0=S
\sqrt{\frac{2\gamma p_0}{mV_0}}}.
\]
\excheck $\gamma=C_p/C_v$ 是绝热指数；
两侧气体都贡献回复力，系数2不可漏。

\oscproblem{63}{162}
图4.22：质量为 $m=21$g 的小球用绝缘细线悬挂在接地无限大水平导体平面上方，
平衡高度 $h=12$cm。
给小球电荷 $q$ 后，振动周期变为原来的 $1/\eta$，
其中 $\eta=2.0$（按书末答案与“changed $\eta$ times”原版语境）。
求 $|q|$。
\exsolution
导体平面镜像等效为距小球 $2h$ 的相反电荷。
球在固定高度附近所受向下静电吸引力
$F_e=q^2/(16\pi\varepsilon_0h^2)$。
在摆动小角度时，球相对平衡点竖直位置变化是二阶小量，
因此可视为有效重力
$g_{\rm eff}=g+F_e/m$，
小振动角频率之比
$\omega_q/\omega_0=\sqrt{g_{\rm eff}/g}=\eta$。
于是
\[
\eta^2=1+\frac{q^2}{16\pi\varepsilon_0h^2mg},
\quad
\boxed{|q|=4h\sqrt{\pi\varepsilon_0mg(\eta^2-1)}
\simeq2.0\,\mathrm{\mu C}}.
\]
\excheck 由于导体吸引加强向下回复力，带电后\emph{周期缩短}；
原版英文仅写“changed $\eta$ times”，
此方向由镜像电力及书后答案共同确认。

\oscproblem{64}{162--163}
小磁针在磁场 $\bm B$ 中绕垂直于磁感应强度的轴作小振动。
改变磁场后周期减小为原来的
$1/\eta$，其中 $\eta=5.0$，
忽略阻尼。
求磁感应强度改变倍数和方向。
\exsolution
磁针磁矩为 $p_m$，
转动惯量 $J$。
小偏角回复力矩
$\tau=-p_mB\theta$，
故
$\omega_0^2=p_mB/J$，
周期
$T=2\pi\sqrt{J/(p_mB)}\propto B^{-1/2}$。
令 $T_f/T_i=1/\eta$，
得到
$\sqrt{B_i/B_f}=1/\eta$，
\[
\boxed{B_f=\eta^2B_i=25B_i}.
\]
\excheck 磁场应\emph{增强}25倍，而不是减小；
摆针本身磁矩与转动惯量保持不变。

\oscproblem{65}{163}
图4.23：两根水平平行导轨间距为 $\ell$，
左端经自感为 $L$ 的理想线圈连接，
右端有质量 $m$ 的活动导杆。
均匀竖直磁场为 $B$，回路总电阻忽略。
$t=0$ 赋予杆向右初速度 $v_0$，
杆初位移为零。
求 $x(t)$。
\exsolution
回路总磁链守恒
$LI+B\ell x=\mathrm{const}$。
以初始电流零、位移零取常数零，
得
$I=-(B\ell/L)x$。
活动杆受安培力
$F_x=IB\ell=-(B^2\ell^2/L)x$。
故
\[
m\ddot x+\frac{B^2\ell^2}{L}x=0,
\qquad
\omega_0=\frac{B\ell}{\sqrt{mL}}.
\]
用初值 $x(0)=0,\ \dot x(0)=v_0$，
\[
\boxed{x(t)=\frac{v_0}{\omega_0}\sin\omega_0t}.
\]
\excheck 无电阻意味着机械能与电感磁场能交换，
没有磁阻尼耗散；轨迹为不衰减正弦振动。
